% Generated draft: requires mathematical and source review.
\documentclass[11pt]{article}
\usepackage[margin=25mm]{geometry}
\usepackage{fontspec}
\setmainfont{Latin Modern Roman}
\usepackage{amsmath,amssymb,mathtools,unicode-math}
\setmathfont{Latin Modern Math}
\usepackage{enumitem,needspace,xurl,hyperref}
\setlength{\parindent}{0pt}
\setlength{\parskip}{5pt}
\setlength{\emergencystretch}{3em}
\newcommand{\N}{\mathbb{N}}
\newcommand{\Z}{\mathbb{Z}}
\newcommand{\Q}{\mathbb{Q}}
\newcommand{\R}{\mathbb{R}}
\newcommand{\C}{\mathbb{C}}
\newcommand{\F}{\mathbb{F}}
\newcommand{\A}{\mathbb{A}}
\newcommand{\PP}{\mathbb{P}}
\newcommand{\E}{\mathbb{E}}
\newcommand{\eps}{\varepsilon}
\newcommand{\dd}{\,\mathrm{d}}
\begin{document}
\section*{UnsolvedMath ID: 30000860}
\textbf{Source problem code:} catalogue entry
\subsection{Definitions and mathematical statement}
The phrase ``the same model'' is clarified by the adjacent local record OWR-1731-004: bond percolation on the square lattice \(\Z^2\), with horizontal and vertical open-edge probabilities \(p_h,p_v\) satisfying \(p_h+p_v=1\). In the conventional independent-edge interpretation, write \(\Pr_{p_h,p_v}\) for the product law and let \(H(a,b)\) be the event that an open nearest-neighbor path in the lattice rectangle \([0,a]\times[0,b]\) joins its left and right sides.

The question asks whether the known rectangular a priori crossing bounds can be improved to give stronger asymptotics for the report's aspect-ratio or shape function \(f(n)\), in particular a nontrivial limit
\[
\lim_{n\to\infty}\frac{n}{f(n)}.
\]
A positive finite limit would be one natural reading of ``nontrivial,'' but the extraction does not define that word or specify a proposed value. More broadly, it asks for a class of shapes admitting both an a priori crossing bound and a conditional extension theorem.

Neither this record nor the adjacent model record defines \(f(n)\), states the existing rectangular estimate, or specifies the conditioning event and conclusion of the extension theorem. Thus no particular threshold definition, rectangle dimension convention for \(f(n)\), rate, lower-bound constant, or conditional probability inequality is inserted here. It is also unclear whether parameters are fixed and constants may depend on them, or whether uniformity over anisotropies is intended; the nondegenerate range \(0<p_h,p_v<1\) is not explicitly stated. These data are required for a checkable asymptotic or crossing theorem.
\subsection{Short English statement}
Can crossing estimates in critical anisotropic square-lattice percolation yield a nontrivial limit of \(n/f(n)\), and can suitable shapes support both an a priori bound and a conditional extension result? The function and estimates are not defined in the extraction.
\subsection{Sources}
\begin{itemize}
\item Local UnsolvedMath registry, permanent ID \(30000860\), source code catalogue entry, ``Crossing Bounds and Shape Functions in Anisotropic Percolation.'' No individual proposer is supplied.
\item Adjacent local registry record, permanent ID \(30000859\), OWR-1731-004, ``RSW Theory for Anisotropic Square-Lattice Percolation,'' supplies the probabilities \(p_h+p_v=1\) and the event \(H(a,b)\) referred to by ``the same model.'' It does not supply a definition of \(f(n)\).
\item \emph{Arbeitsgemeinschaft: Percolation} (2008), Oberwolfach Reports, DOI 10.4171/owr/2007/48, pp. 2851--2892: \url{https://doi.org/10.4171/owr/2007/48}. This is the sole supplied external reference.
\item Supplied literature triage, checked 2026-08-21: status ``open.'' It describes this as a deliberately broad companion shape-function question and reports no identified general resolution. This dated assessment derives from supplied OWR curation; time-sensitive status and cited claims require independent verification.
\end{itemize}
\end{document}
